\section{Deriving the probability distribution over number of prey eaten }

To perform Bayesian inference for our model, we need to calculate the likelihood, which is composed of individual (assumed independent) contributions from each experimental replicate.

Each replicate begins with an initial number of prey, denoted $N_0$. For a given $N_0$, we seek the probability that at time $t=1 \text{ hour}$ (the end of the experiment), there are $N$ prey remaining, where $0 \leq N \leq N_0$. To compute this, we model the time evolution of the system using the \emph{chemical master equation} (CME), which describes how the probability of each system state changes over time.

We represent the predator state using a binary variable:

\begin{equation}
s =
\begin{cases}
1 & \text{if the predator is in the active preying (AP) state}, \\
0 & \text{if the predator is in the between-preying (BP) state}.
\end{cases}
\end{equation}

The full system state is given by the vector $(N, s)$, where $N$ is the number of prey remaining. Let $p_{N,s}(t)$ denote the probability that the system is in state $(N,s)$ at time $t$. The CME defines a system of ordinary differential equations governing the evolution of $p_{N,s}(t)$. For $0 \leq N \leq N_0$, the equations take the form:
\begin{align}
\label{eq:ap}
\frac{d}{dt} p_{N,1}(t) &= -a N \, p_{N,1}(t) + r \, p_{N,0}(t), \\
\label{eq:bp}
\frac{d}{dt} p_{N,0}(t) &= a (N+1)\, p_{N+1,1}(t) - r \, p_{N,0}(t),
\end{align}

where $a$ is the prey capture rate and $r$ is the prey handling rate. For $N=0$, eq. \eqref{eq:ap} is modified by removing the outflow term from this state; for $N = N_0$, eq. \eqref{eq:bp} is modified by removing the inflow term from the $(N+1)$th state.

At time $t = 0$, the system starts with all prey present and the predator actively preying:

\[
p_{N_0,1}(0) = 1, \quad p_{N,s}(0) = 0 \text{ otherwise}.
\]

Since the system begins in state $(N_0,1)$, the state $(N_0,0)$ is never reached, because entering the BP state requires consuming at least one prey item.

As $t \to \infty$, the system approaches the absorbing state $(0,1)$ corresponding to no prey remaining and the predator still searching -- a hungry fish left in an empty bucket.

\subsection*{Matrix formulation}
The system of equations can be solved analytically using Fourier transforms, although the length of the solutions, in terms of symbols, grows exponentially with $N_0$. For efficiency, instead, we solve the equations numerically. To solve the system, we stack the state probabilities into a vector:
\begin{equation}
\mathbf{p}(t) =
\begin{pmatrix}
p_{N_0,0}(t), & p_{N_0,1}(t), & \cdots, & p_{1,0}(t), & p_{1,1}(t), & p_{0,0}(t), & p_{0,1}(t)
\end{pmatrix}.
\end{equation}

The CME can then be written as a linear dynamical system:
\[
\frac{d}{dt} \mathbf{p}(t) = \boldsymbol{A} \, \mathbf{p}(t),
\]
where $\boldsymbol{A}$ is a transition-rate matrix defined by the reaction structure:

\begin{equation}
\boldsymbol{A} =
\begin{pmatrix}
-r & 0 & 0 & 0 & \cdots & 0 & 0 & 0\\
r & -aN_0 & 0 & 0 & \cdots & 0 & 0 & 0\\
0 & aN_0 & -r & 0 & \cdots & 0 & 0 & 0\\
0 & 0 & r & -a(N_0-1) & \ddots & \vdots & \vdots & \vdots\\
\vdots & \vdots & 0 & a(N_0-1) & \ddots & 0 & 0 & 0\\
\vdots & \vdots & \vdots & \ddots & \ddots & -a & 0 & 0\\
0 & 0 & \cdots & 0 & 0 & a & -r & 0\\
0 & 0 & \cdots & 0 & 0 & 0 & r & 0
\end{pmatrix}.,
\end{equation}

The solution is:
\[
\mathbf{p}(t) = \exp(\boldsymbol{A} t)\, \mathbf{p}(0),
\]
where $\exp(\cdot)$ denotes the matrix exponential.

To obtain the probability distribution over the number of prey remaining, we marginalise over the predator state:
\[
P(N,t) = p_{N,0}(t) + p_{N,1}(t).
\]

Here, $P(N,t)$ is the probability that $N$ prey remain at time $t$.

\subsection*{Likelihood construction}

For a given experimental replicate $i$, suppose $N_i$ prey are observed at time $t=1 \text{ hour}$. The likelihood contribution is:
\[
\mathcal{L}_i = P(N_i, t=1),
\]
where $P(N_i,t=1)$ is computed via the solution above, using parameters $a$ and $r$ (which depend on covariates through eqs. \eqref{eqn:eqn3} \& \eqref{eqn:eqn4} in the main text).

Assuming independence across replicates, the full likelihood is:
\[
\mathcal{L} = \prod_i P(N_i, t=1).
\]

Because each replicate may have different covariates and initial prey counts $N_0$, the probability $P(N,t)$ must be computed separately for each replicate.

\section{A limiting case: voracious foragers ($a \to \infty$)}
\label{app:voracious}

The likelihood above under-predicts the fraction of fish that consume \emph{all} of their prey. This motivates a limiting version of the model in which prey capture is effectively instantaneous, so that a fish is constrained only by the handling (BP) process. Formally, we take $a \to \infty$ in eq.\ \eqref{eqn:eqn1} of the main text. In this limit the sojourn time in the AP state with $N \geq 1$ prey remaining, which is exponentially distributed with rate $aN$, tends to zero: whenever the predator is in the AP state and at least one prey remains, it captures a prey item and transitions to the BP state \emph{immediately} (with probability one).

Consider a replicate that begins with $N_0 \geq 1$ prey. The dynamics are then as follows. At $t = 0$ the predator is in the AP state, so it instantaneously captures its first prey and enters the BP state. It remains there for a handling period that is exponentially distributed with rate $r$; upon completion it re-enters the AP state and, provided prey remain, instantaneously captures another and returns to BP; and so on. Every capture after the first is therefore preceded by exactly one $\mathrm{Exponential}(r)$ handling period, while the instantaneous captures consume no time. The completion times of successive handling periods thus form a homogeneous Poisson process of rate $r$, and the number of such completions during the trial of unit length,
\[
K \sim \mathrm{Poisson}(r),
\]
counts the captures made \emph{after} the initial instantaneous one. The process stops when either the trial ends or the prey are exhausted, so the total number of prey eaten is
\begin{equation}
\label{eq:voracious-count}
Y = \min\!\big(N_0,\; 1 + K\big), \qquad K \sim \mathrm{Poisson}(r).
\end{equation}
This is a shifted (by one), right-truncated (at $N_0$) Poisson distribution. Explicitly, its probability mass function is
\begin{equation}
\label{eq:voracious-pmf}
P_{\mathrm{vor}}(Y = k \mid r, N_0) =
\begin{cases}
\dfrac{e^{-r} r^{\,k-1}}{(k-1)!}, & 1 \leq k \leq N_0 - 1, \\[2ex]
\displaystyle 1 - \sum_{j=0}^{N_0 - 2} \dfrac{e^{-r} r^{\,j}}{j!}, & k = N_0, \\[2ex]
0, & \text{otherwise.}
\end{cases}
\end{equation}
Because $N_0 \geq 1$, a voracious forager always eats at least one prey item ($P_{\mathrm{vor}}(Y = 0) = 0$); the probability mass that a Poisson count would place above $N_0 - 1$ is accumulated onto the full-consumption outcome $Y = N_0$, which is precisely the behaviour the CME model failed to capture.

\subsection*{Mixture likelihood}

We allow each fish to be either an ordinary forager, whose consumption follows the CME distribution $P(N_i, t = 1)$ derived above, or a voracious forager, whose consumption follows eq.\ \eqref{eq:voracious-pmf}. Let $z_i \in \{0, 1\}$ indicate that replicate $i$ is voracious, with $\Pr(z_i = 1) = \pi_i$ and $\mathrm{logit}(\pi_i) = \gamma_0 + \gamma_w \, \mathrm{weight}_i$. Writing $y_i$ for the number of prey eaten by fish $i$ (so $y_i = N_{0,i} - N_i$) and marginalising over the latent indicator $z_i$, the likelihood contribution is
\begin{equation}
\label{eq:mixture-like}
\mathcal{L}_i = (1 - \pi_i)\, P_{\mathrm{CME}}(y_i \mid a_i, r_i, N_{0,i}) \; + \; \pi_i \, P_{\mathrm{vor}}(y_i \mid r_i, N_{0,i}),
\end{equation}
where $P_{\mathrm{CME}}$ is the CME distribution over prey eaten (equivalently $P(N_i, t=1)$ with $N_i = N_{0,i} - y_i$). The full likelihood is the product of eq.\ \eqref{eq:mixture-like} across independent replicates. Setting $\pi_i \equiv 0$ recovers the pure CME model, so the two-state model is nested within the mixture.